\section{Languages}

\subsection{Languages Spoken}

\subsection{Official Languages}

\subsection{Languages}
Considering these language-related factors allows for a
comprehensive evaluation of a country's linguistic landscape,
acknowledging the importance of languages in cultural expression,
communication, education, and social cohesion.
\begin{enumerate}
  \item \textbf{Official Language(s)}: The language(s) designated
  as official at the national level, impacting government,
  education, and legal proceedings.
  
  \item \textbf{Language Diversity}: The number and variety of
  languages spoken within the country, which can affect cultural
  richness and linguistic heritage.
  
  \item \textbf{Language Proficiency}: The level of fluency and
  proficiency in the official language(s) and other widely spoken
  languages, influencing communication and social integration.
  
  \item \textbf{Language Education}: The availability and quality
  of language education programs, promoting multilingualism and
  cultural understanding.
  
  \item \textbf{Language Policy}: The government's approach to
  language use, including support for minority languages and
  language preservation efforts.
  
  \item \textbf{Language Access}: The extent to which language
  barriers are addressed in public services, healthcare, and
  legal systems, ensuring equitable access for all citizens.
  
  \item \textbf{Language in Media and Literature}: The
  representation and promotion of languages in media,
  literature, and online platforms, fostering linguistic
  diversity.
  
\end{enumerate}


